% ================================================================
%  第 7 章  三流形中的极小曲面
% ================================================================
\chapter{三流形中的极小曲面}
\label{chap:7}

{\setlength{\emergencystretch}{2em}%
本章讨论三流形中极小曲面的理论. 我们先说明如何把先前的结论推广到这一
情形 (特别是单调性、强极值原理, 以及关于极小曲面的其他一些基本估计).
下一节证明 Hersch 以及 Yang 与 Yau 的特征值估计. 之后转向 Reilly 的一个
积分公式与 Choi--Wang 关于嵌入极小曲面的特征值估计. 接着证明 Choi 与
Schoen 关于正 Ricci 曲率三流形中嵌入极小曲面的紧性定理; 这一紧性结论的
要点在于: 由 Choi--Wang 以及 Yang 与 Yau 的结果, 这类极小曲面具有一致
的面积界. 再下一节证明 Schoen 与 Yau 的\newterm{positive-mass}{正质量定理}{positive mass theorem}.
最后, 在末节证明 Colding--Minicozzi 关于同伦球面上 Ricci 流熄灭时间的
估计.\par}

% ================================================================
\section{三流形中的极小曲面方程}
\label{sec:7.1}

本节先描述讨论一般 Riemann 三流形中极小曲面所需的改动. 局部图像与
$\R^3$ 中极小曲面的情形非常相似. 事实上, 任意三流形中极小曲面的许多
正则性结果, 都可以借助重新尺度化的论证归约到 $\R^3$ 中相应的整体结论.

以下 $M^n$ 总表示完备的 $n$ 维 Riemann 流形, 其截面曲率有界 (即
$|K_M|\leq k$), 单射半径有正下界 $i_0>0$. 通常取 $n=3$.

为建立联系曲率与面积所必需的局部结果, 我们需要回顾一些预备的几何事实.
标准的 Hessian 比较定理蕴含:

\begin{lemma}
\label{lem:7.1}
对 $r<\min\{i_0,1/\sqrt{k}\}$ 及任意满足 $|X|=1$ 的向量 $X$,
\begin{equation}
\Bigl|\Hess_r(X,X)-\frac1r\bigl\langle X-\langle X,Dr\rangle Dr,
X-\langle X,Dr\rangle Dr\bigr\rangle\Bigr|\leq\sqrt{k}.
\label{eq:7.1}
\end{equation}
\end{lemma}

\begin{proof}
距离函数的 Hessian 在径向为零. 由 Hessian 比较定理, Hessian 其余的特征
值分别以 $\sqrt{k}\coth\sqrt{k}r$ 与 $\sqrt{k}\cot\sqrt{k}r$ 为上、下界.
由此, 用初等不等式即得估计 \eqref{eq:7.1}.
\end{proof}

设 $x\in\Sigma^2\subset M^3$ 而 $\Sigma$ 极小. 当 $M$ 为 $\R^n$ 时,
$\Sigma$ 的极小性蕴含
\begin{equation}
\Delta r^2=4.
\label{eq:7.2}
\end{equation}
这一事实是证明极小曲面的单调性公式 (即命题~\ref{prop:1.12}) 的关键. 在一般情形,
利用极小性与 \eqref{eq:7.1} 可得
\begin{equation}
\bigl|\Delta r^2-4\bigr|\leq4\sqrt{k}\,r
\label{eq:7.3}
\end{equation}
对 $r<\min\{i_0,1/\sqrt{k}\}$ 成立. 应用余面积公式 (即 \eqref{eq:1.59}) 并用
Stokes 定理, 给出
\begin{equation}
s\frac{d}{ds}\Area(B_s\cap\Sigma)
=\int_{\partial B_s\cap\Sigma}\frac{r}{|\nabla r|}
\geq\frac12\int_{B_s\cap\Sigma}\Delta r^2.
\label{eq:7.4}
\end{equation}
把 \eqref{eq:7.3} 与 \eqref{eq:7.4} 结合起来, 可知对
$s<\min\{i_0,1/\sqrt{k}\}$ 有
\begin{equation}
\frac{d}{ds}\Bigl(e^{2\sqrt{k}s}s^{-2}\Area(B_s\cap\Sigma)\Bigr)\geq0.
\label{eq:7.5}
\end{equation}

我们可以像命题~\ref{prop:1.15} 那样论证, 得到三流形中极小曲面的平均值不等式.

当 $M=\R^3$ 时, 极小曲面的 Gauss 映射是共形的. 更一般地, Gauss 方程与
极小性共同蕴含
\begin{equation}
K_{\Sigma}=K_M-\frac12|A|^2,
\label{eq:7.6}
\end{equation}
从而
\begin{equation}
|A|^2\leq2k-2K_{\Sigma}.
\label{eq:7.7}
\end{equation}
特别地, \eqref{eq:7.6} 蕴含 Gauss 映射是拟共形的. 称映射
$F:(M_1,g_1)\to(M_2,g_2)$ 为\newterm{quasi-conformal}{拟共形}{quasi-conformal},
若存在常数 $\Lambda<\infty$, 使得对一切 $x\in M_1$, $F^*g_2$ 的最大特征值
与最小特征值之比不超过 $\Lambda$. 注意共形映射必然是拟共形的 (取
$\Lambda=1$).

本节余下部分在局部坐标 $(x_1,x_2,x_3)$ 与 $M$ 上的度量 $g_{ij}$ 下工作.
置 $e_i$ 为向量场 $\frac{\partial}{\partial x_i}$, 于是
$\langle e_i,e_j\rangle=g_{ij}$. 用 $\Gamma_{ij}^n$ 表示相应 Riemann
联络的 Christoffel 符号.

设 $u:\Omega\subset\R^2\to\R$ 是 $C^2$ 函数, 考虑 $u$ 的图
\begin{equation}
\Sigma=\Graph_u
=\bigl\{(x_1,x_2,u(x_1,x_2))\mid(x_1,x_2)\in\Omega\bigr\}.
\label{eq:7.8}
\end{equation}
我们先在这些坐标下导出 $\Sigma$ 的极小曲面方程. 为此需要把 $\Sigma$ 的
平均曲率用 $u$ 表示出来.

对 $i=1,2$, 定义向量场 $E_i$ 与线性映射 $T_i^n$ 为
\begin{equation}
E_i=e_i+u_{x_i}e_3\equiv T_i^ne_n,
\label{eq:7.9}
\end{equation}
于是 $(E_1,E_2)$ 构成 $\Sigma$ 切空间的一组基. 记 $h_{ij}=\langle E_i,E_j\rangle$
为 $\Sigma$ 上的诱导度量. 方便起见, 定义 $W$ 为
\begin{equation}
W^2=1+g^{ij}u_{x_i}u_{x_j}.
\label{eq:7.10}
\end{equation}
若 $N$ 表示 $\Sigma$ 的向上的单位法向量, 则 $\langle N,E_i\rangle=0$. 故对
$i=1,2$ 有
\begin{equation}
\langle N,e_i\rangle=-\frac{u_{x_i}}{W}
\label{eq:7.11}
\end{equation}
以及
\begin{equation}
\langle N,e_3\rangle=\frac1W.
\label{eq:7.12}
\end{equation}
平均曲率由
\begin{equation}
h^{ij}\bigl\langle N,\nabla_{E_i}E_j\bigr\rangle
\label{eq:7.13}
\end{equation}
给出. 利用 \eqref{eq:7.9} 算得
\begin{equation}
h^{ij}=g^{ij}-W^{-2}g^{in}g^{jl}u_{x_n}u_{x_l}
\label{eq:7.14}
\end{equation}
以及
\begin{equation}
\begin{aligned}
\nabla_{E_i}E_j
&=T_i^n\nabla_{e_n}T_j^le_l
=T_i^nT_j^l\nabla_{e_n}e_l+(T_i^n\nabla_{e_n}T_j^l)e_l\\
&=T_i^nT_j^l\Gamma_{nl}^me_m+E_i(T_j^l)e_l
=T_i^nT_j^l\Gamma_{nl}^me_m+u_{x_ix_j}e_3.
\end{aligned}
\label{eq:7.15}
\end{equation}
结合 \eqref{eq:7.11}、\eqref{eq:7.12} 与 \eqref{eq:7.15}, 得到
\begin{equation}
\langle N,\nabla_{E_i}E_j\rangle
=-\frac1W\sum_{m=1}^2\Bigl(u_{x_m}T_i^nT_j^l\Gamma_{nl}^m\Bigr)
+\frac1W\Bigl(u_{x_ix_j}+T_i^nT_j^l\Gamma_{nl}^3\Bigr).
\label{eq:7.16}
\end{equation}
用 $W$ 与 $h^{ij}$ 乘遍全式, 即得极小曲面方程
\begin{equation}
0=h^{ij}\Bigl(u_{x_ix_j}+T_i^nT_j^l\Gamma_{nl}^3\Bigr)
-\sum_{m=1}^2\Bigl(u_{x_m}h^{ij}T_i^nT_j^l\Gamma_{nl}^m\Bigr).
\label{eq:7.17}
\end{equation}
代入 $T_i^n=\delta_i^n+u_{x_i}\delta_3^n$ 的定义, 可写
\begin{equation}
\begin{aligned}
T_i^nT_j^l\Gamma_{nl}^m
&=(\delta_i^n+u_{x_i}\delta_3^n)(\delta_j^l+u_{x_j}\delta_3^l)
\Gamma_{nl}^m\\
&=\Gamma_{ij}^m+u_{x_i}\Gamma_{3j}^m+u_{x_j}\Gamma_{i3}^m
+u_{x_i}u_{x_j}\Gamma_{33}^m.
\end{aligned}
\label{eq:7.18}
\end{equation}
把 \eqref{eq:7.18} 代入 \eqref{eq:7.17},
\begin{equation}
\begin{aligned}
0={}&h^{ij}\Bigl(u_{x_ix_j}+\Gamma_{ij}^3+u_{x_i}\Gamma_{3j}^3
+u_{x_j}\Gamma_{i3}^3+u_{x_i}u_{x_j}\Gamma_{33}^3\Bigr)\\
&-\sum_{m=1}^2u_{x_m}h^{ij}\Bigl(\Gamma_{ij}^m+u_{x_i}\Gamma_{3j}^m
+u_{x_j}\Gamma_{i3}^m+u_{x_i}u_{x_j}\Gamma_{33}^m\Bigr).
\end{aligned}
\label{eq:7.19}
\end{equation}

接着用它定义函数 $F:\R^9\to\R$ 为
\begin{equation}
\begin{aligned}
F(x_1,x_2,u,p_1,p_2,p_{ij})
={}&h^{ij}\Bigl(p_{ij}+\Gamma_{ij}^3+p_i\Gamma_{3j}^3
+p_j\Gamma_{i3}^3+p_ip_j\Gamma_{33}^3\Bigr)\\
&-\sum_{m=1}^2p_mh^{ij}\Bigl(\Gamma_{ij}^m+p_i\Gamma_{3j}^m
+p_j\Gamma_{i3}^m+p_ip_j\Gamma_{33}^m\Bigr),
\end{aligned}
\label{eq:7.20}
\end{equation}
其中 $\Gamma_{ij}^n=\Gamma_{ij}^n(x_1,x_2,u)$, 而 $h^{ij}$ 是
\[
h_{ij}=g_{ij}(x_1,x_2,u)+p_ig_{j3}(x_1,x_2,u)
+p_jg_{3i}(x_1,x_2,u)+p_ip_jg_{33}(x_1,x_2,u)
\]
的逆矩阵. 按构造, 若 $u$ 是极小曲面方程的解, 则
\begin{equation}
F(x_1,x_2,u,u_{x_1},u_{x_2},u_{x_ix_j})=0.
\label{eq:7.21}
\end{equation}
若 $\nabla u$ 有界, 则 \eqref{eq:7.21} 是一致椭圆的. 于是当
$|\nabla u|$ 有界时, 可用 $|u|$ 的最大值给出 $u$ 的 $C^{2,\alpha}$ 估计.

我们还要给出 \eqref{eq:7.21} 之形式的另外两个推论: 先是一个可去奇点结果,
然后是两个极小曲面相交的局部描述.

下面的可去奇点结果远非最优. 事实上, 对极小曲面方程而言, 可去奇点方面的
结论比线性椭圆方程的情形强得多 (例如可参看 \cite{Os2} 第 10 章). 不过
这一结果足以应付我们后面的应用.

\begin{lemma}
\label{lem:7.2}
设 $u$ 是 $B_1\subset\R^2$ 上的 $C^1$ 函数, 且
\begin{equation}
|u|+|\nabla u|\leq C\quad\text{于}\ B_1.
\label{eq:7.22}
\end{equation}
若 $u$ 是极小曲面方程在 $C^2(B_1\setminus\{0\})$ 中的解, 则 $u$ 是整个
$B_1$ 上的光滑解.
\end{lemma}

\begin{proof}
假设 \eqref{eq:7.22} 蕴含极小曲面方程一致椭圆且系数有界. 把该方程写成
散度形式, 则一致梯度界 \eqref{eq:7.22} 与 $u$ 在 $0$ 之外是 $C^2$ 解这一
事实共同蕴含 $u$ 是整个 $B_1$ 上的 $W^{1,2}$ 弱解. 椭圆估计
(\cite{GiTr} 定理 8.8) 蕴含 $u\in H^2(B_1)$, 从而 $|\nabla u|\in W^{1,2}(B_1)$.
\end{proof}

现在设 $u$ 与 $v$ 都是 \eqref{eq:7.21} 的光滑解. 令 $w=v-u$, 并对
$0\leq s\leq1$ 定义
\begin{equation}
G(x_1,x_2,s)=F\bigl(x_1,x_2,u+sw,(u+sw)_{x_1},(u+sw)_{x_2},
(u+sw)_{x_ix_j}\bigr).
\label{eq:7.23}
\end{equation}
由微积分基本定理与链式法则,
\begin{equation}
\begin{aligned}
0&=G(x_1,x_2,1)-G(x_1,x_2,0)
=\int_{s=0}^1\frac{d}{ds}G(x_1,x_2,s)\,ds\\
&=\int_{s=0}^1F_uw+F_{p_i}w_{x_i}+F_{p_{ij}}w_{x_ix_j}\,ds,
\end{aligned}
\label{eq:7.24}
\end{equation}
其中 $F$ 的每个偏导数都在
\begin{equation}
\bigl(x_1,x_2,u+sw,u_{x_1}+sw_{x_1},u_{x_2}+sw_{x_2},
u_{x_ix_j}+sw_{x_ix_j}\bigr)
\label{eq:7.25}
\end{equation}
处取值. 特别地, \eqref{eq:7.24} 表明 $w$ 自身满足一个偏微分方程, 其系数
是 $F$ 的偏导数的积分.

这一对极小曲面方程的描述使我们能应用 Bers 定理 (即定理~\ref{thm:6.12}) 来刻画三
流形中极小曲面的相交.

\begin{theorem}[极小曲面相交的局部描述]
\label{thm:7.3}
设 $\Sigma_1^2,\Sigma_2^2\subset M^3$ 是光滑连通浸入极小曲面, 且它们不在
任何开集上重合. 则 $\Sigma_1$ 与 $\Sigma_2$ 横截相交, 只在某个孤立点集
$\mathcal{D}$ 处例外. 给定 $y\in\mathcal{D}$, 存在整数 $d\geq2$ 与 $y$ 的
一个邻域 $U$, 使得在 $U$ 内的相交由 $2d$ 条在 $y$ 处相交的嵌入弧组成.
\end{theorem}

\begin{proof}
不妨设 $y\in\mathcal{D}$ 是非横截相交点. 选取坐标使 $y=(0,0,0)$. 可设
$\Sigma_1$ 与 $\Sigma_2$ 都是 $(0,0,0)$ 处公共切平面上的函数 $u$ 与 $v$ 的
图. 于是 $u,v,u_{x_i},v_{x_i}$ 在 $(0,0)$ 处全为零. 定义 $w=u-v$.

先证明 $w$ 满足一个系数光滑的一致椭圆微分方程
\begin{equation}
a_{ij}w_{x_ix_j}+b_iw_{x_i}+cw=0.
\label{eq:7.26}
\end{equation}
由于只关心局部性质, 可设 $u,v$ 及其导数一致有界. 由此立即得到
\eqref{eq:7.24} 中的系数光滑有界, 且 $a_{ij}$ 对称. 为验证椭圆性, 需要
计算
\begin{equation}
a_{ij}=\int_{s=0}^1F_{p_{ij}}\bigl(x_1,x_2,u+sw,(u+sw)_{x_1},
(u+sw)_{x_2},(u+sw)_{x_ix_j}\bigr)\,ds.
\label{eq:7.27}
\end{equation}
由 (微积分中的) 中值定理, 对每个 $(x_1,x_2)$ 存在某个 $0\leq t\leq1$,
使得若 $f=u+tw$, 则
\begin{equation}
a_{ij}(x_1,x_2)=F_{p_{ij}}\bigl(x_1,x_2,f,f_{x_1},f_{x_2},
f_{x_ix_j}\bigr).
\label{eq:7.28}
\end{equation}
容易看出 $F_{p_{ij}}=h^{ij}$, 其中 $h^{ij}$ 是
\begin{equation}
h_{ij}=g_{ij}+f_{x_i}g_{j3}+f_{x_j}g_{3i}+f_{x_i}f_{x_j}g_{33}
\label{eq:7.29}
\end{equation}
的逆矩阵, 而式中
\[
g_{nl}=g_{nl}(x_1,x_2,f(x_1,x_2)).
\]
显然 $h^{ij}$ 一致椭圆当且仅当 $h_{ij}$ 一致椭圆. 最后, 由于 $g_{ij}$ 按
定义一致椭圆, 且可设 $f_{x_i}$ 任意小, 故可断定 $a_{ij}$ 在某个邻域
$U\subset\R^2$ 内一致椭圆. 既然 $w$ 满足 \eqref{eq:7.26}, 定理~\ref{thm:6.12} 蕴含:
或者 $w$ 恒为零, 或者存在线性变换 $T$, 使得用 $T$ 旋转之后可写
\begin{equation}
w(x_1,x_2)=p_d(x_1,x_2)+q(x_1,x_2),
\label{eq:7.30}
\end{equation}
其中 $p_d$ 是次数 $d<\infty$ 的齐次调和多项式, 且
\begin{equation}
|q(x)|+|x||\nabla q(x)|+\cdots+|x|^d|\nabla^dq(x)|\leq C|x|^{d+1}.
\label{eq:7.31}
\end{equation}
特别地, 这表明非横截相交点是孤立的. 最后, 由引理~\ref{lem:6.13} 的一个小变形
立即得到 $u(U)\cap v(U)$ 同胚于在 $x$ 处相交的 $2d$ 条弧.
\end{proof}

定理~\ref{thm:7.3} 及其在 $\Sigma_i$ 带分支情形下的变体 (参看定理~\ref{thm:6.18}), 可用来把
第 4 章的结论推广到三流形 $M^3$ 中的极小曲面.

作为定理~\ref{thm:7.3} 的一个应用, 我们将证明 J. Hass 关于带极小叶状结构的三流形
中极小曲面的一个结果.

在解释这一结果之前, 我们需要层状结构的定义 (例如可参看 W. P. Thurston
\cite{Th} 第 8 章, D. Gabai \cite{Ga}, 或 D. Calegari \cite{Ci}).

\begin{definition}[层状结构]
\label{def:7.4}
$M^3$ 的一个余维一\newterm{lamination}{层状结构}{lamination}是一族互不相交的
光滑连通曲面 (称为\emph{叶}), 使得 $\bigcup_{\Lambda\in\mathcal{L}}\Lambda$
是闭集. 此外, 对每个 $x\in M$ 存在 $x$ 的开邻域 $U$ 与局部坐标卡
$(U,\Phi)$, 其中 $\Phi(U)\subset\R^3$, 使得在这些坐标下 $\mathcal{L}$ 中
的叶经过该坐标卡时呈形如 $(\R^2\times\{t\})\cap\Phi(U)$ 的薄片.
\end{definition}

\newterm{foliation}{叶状结构}{foliation}是叶的并集等于整个 $M$ 的层状结构.
若叶是可定向的, 则称该叶状结构是可定向的. 若所有叶都是极小子流形, 则称
该叶状结构是极小的.

\begin{theorem}[J. Hass, \cite{Ha}]
\label{thm:7.5}
若 $M^3$ 具有一个无紧叶的可定向余维一极小叶状结构, 则 $M$ 不含浸入的
极小 $S^2$.
\end{theorem}

\begin{proof}
反设存在浸入极小球面 $\Sigma\subset M$, 导出矛盾. 由于没有一片叶是紧的,
$\Sigma$ 不会含于某片叶中. 因此, 若 $p\in\Sigma$, $p\in\Lambda\in\mathcal{L}$
且 $T_p\Sigma=T_p\Lambda$, 则定理~\ref{thm:7.3} 蕴含存在 $p$ 的邻域 $U$, 使得二者在
$U\setminus\{p\}$ 内横截相交. 特别地, 叶状结构限制到 $\Sigma$ 上给出一个
奇异的叶状结构, 其向量场在 (非横截相交的) 孤立点处为零. 应用定理~\ref{thm:7.3}
可知, 该向量场在每个零点处的指标为负. Poincaré--Hopf 公式蕴含这些指标
之和等于 $\Sigma$ 的 Euler 示性数 (它等于二). 这就导出矛盾, 定理得证.
\end{proof}

显然, 这一论证的种种变体可以更一般地应用, 给出类似的结果. 此外, 对
$M^3$ 的某些拓扑假设可使我们证明极小浸入球面必定存在. 第 4 章与第 5 章
叙述了分支极小浸入的一些存在性结果, 而第 6 章说明了在若干情形下如何
排除分支点. 用完全不同的方法, F. Smith 在 \cite{Sm} 中证明了 $S^3$
(带任意度量) 总含有嵌入极小 $S^2$; 另可参看 \cite{CD}. 在这些情形下,
定理~\ref{thm:7.5} 蕴含 $M$ 的每个极小叶状结构都有紧叶 (更多这方面的内容可参看
\cite{Ha}).

最后我们指出, 极小叶状结构因拓扑上的理由而有趣. 1979 年, Dennis Sullivan
\cite{Du2} 证明了: 一个余维一叶状结构是绷紧的, 当且仅当存在使每片叶都
极小的 Riemann 度量. 称三流形的余维一叶状结构是\newterm{taut}{绷紧的}{taut}
的, 若存在一条与每片叶都相交的单一横截圆周, 其中``横截圆周''指一条处处与
叶状结构的切场横截的闭曲线.

% ================================================================
%  7.2  Hersch 与 Yang--Yau 的定理
% ================================================================
\section{Hersch 与 Yang--Yau 的定理}
\label{sec:7.2}

1970 年, J. Hersch 证明了: $\bS^2$ 上的圆度量使第一特征值与面积之积达到
最大. 关键在于一个巧妙的平衡论证: 它使 Hersch 能把到圆球面的共形映射的
坐标函数用作 Rayleigh 商的试验函数. 这一想法由 P. Yang 与 S. T. Yau 在
\cite{YgYa} 中推广到了一般 Riemann 曲面.

\begin{theorem}[Hersch, \cite{Her}]
\label{thm:7.6}
若 $\Sigma^2$ 是带 Riemann 度量的拓扑 $\bS^2$, 则
\begin{equation}
\lambda_1(\Sigma)\Area(\Sigma)
\leq\lambda_1(\bS^2)\Area(\bS^2)=8\pi,
\label{eq:7.32}
\end{equation}
等号成立当且仅当该度量是圆度量.
\end{theorem}

\begin{proof}
由单值化定理 (参看引理~\ref{lem:5.22}), 存在共形微分同胚
\[
\Phi:\Sigma\to\bS^2\subset\R^3.
\]
对每个 $i=1,2,3$, 置 $\phi_i=x_i\circ\Phi$. 给定 $x\in\bS^2$, 令
$\pi_x:\bS^2\setminus\{x\}\to\C$ 为球极投影, 并令
\begin{equation}
\psi_{x,t}(y)=\pi_x^{-1}\bigl(t(\pi_x(y))\bigr),
\label{eq:7.33}
\end{equation}
则对每个 $t,x$, 它都可延拓为 $\bS^2$ 上的共形映射. 定义
$\Psi:\bS^2\times(0,1)\to G$, 其中 $G$ 是 $\bS^2$ 的共形变换群, 为
$\Psi(x,t)=\psi_{x,1/(1-t)}$. 由于对每个 $x\in\bS^2$ 有
$\Psi(x,0)=\mathrm{id}_{\bS^2}$, 可把 $\Psi$ 视为
$B_1(0)=\bS^2\times(0,1)/(x,0)\equiv(y,0)$ 上的连续映射. 置
\begin{equation}
\mathcal{A}\bigl(\Psi(x,t)\bigr)
=\frac1{\Area(\Sigma)}
\Bigl(\int_{\Sigma}x_i\circ\Psi(x,t)\circ\Phi\Bigr)_{i=1,2,3}
\label{eq:7.34}
\end{equation}
于是
\begin{equation}
\mathcal{A}:B_1(0)\to\R^3\quad\text{且}\quad
\lim_{(y,t)\to(x,1)}\mathcal{A}\bigl(\Psi(y,t)\bigr)=x.
\label{eq:7.35}
\end{equation}
特别地, $\mathcal{A}$ 可延拓到 $\partial B_1$ 上成为度为 $1$ 的映射. 因此
由初等拓扑 (必要时把 $\Phi$ 换成 $\psi\circ\Phi$) 可设对每个 $i$,
\begin{equation}
\int_{\Sigma}\phi_i=0;
\label{eq:7.36}
\end{equation}
也就是说, 每个 $\phi_i$ 都与常函数正交. 由第一特征值的通常变分刻画, 对
每个 $i$ 有
\begin{equation}
\int_{\Sigma}|\nabla\phi_i|^2\geq\lambda_1(\Sigma)\int_{\Sigma}\phi_i^2.
\label{eq:7.37}
\end{equation}
对 $i$ 求和, 并利用 $\Phi(\Sigma)\subset\bS^2$ (故
$\sum_{i=1}^3\phi_i^2=1$), 得
\begin{equation}
\sum_{i=1}^3\int_{\Sigma}|\nabla\phi_i|^2
\geq\lambda_1(\Sigma)\Area(\Sigma).
\label{eq:7.38}
\end{equation}
现在, 显然由于 $\Phi$ 共形 (故保持能量), 且每个 $x_i$ 都是
$\bS^2\subset\R^3$ 上特征值为 $\lambda_1(\bS^2)$ 的特征函数, 我们得到
\begin{equation}
\int_{\Sigma}|\nabla\phi_i|^2
=\int_{\bS^2}|\nabla x_i|^2
=\lambda_1(\bS^2)\int_{\bS^2}x_i^2.
\label{eq:7.39}
\end{equation}
把 \eqref{eq:7.38} 与 \eqref{eq:7.39} 结合, 得
\begin{equation}
\begin{aligned}
\lambda_1(\bS^2)\Area(\bS^2)
&=\sum_{i=1}^3\lambda_1(\bS^2)\int_{\bS^2}x_i^2
=\sum_{i=1}^3\int_{\bS^2}|\nabla x_i|^2\\
&=\sum_{i=1}^3\int_{\Sigma}|\nabla\phi_i|^2
\geq\lambda_1(\Sigma)\Area(\Sigma).
\end{aligned}
\label{eq:7.40}
\end{equation}

在 \eqref{eq:7.32} 取等号的情形, 我们看到 \eqref{eq:7.37} 对每个 $i$ 都
取等号, 故必有
\begin{equation}
\Delta\phi_i+\lambda_1(\Sigma)\phi_i=0.
\label{eq:7.41}
\end{equation}
再次利用 $\sum_{i=1}^3\phi_i^2=1$, 得
\begin{equation}
\begin{aligned}
0&=\Delta\sum_{i=1}^3\phi_i^2
=2\sum_{i=1}^3\phi_i\Delta\phi_i+2\sum_{i=1}^3|\nabla\phi_i|^2\\
&=-2\lambda_1(\Sigma)\sum_{i=1}^3\phi_i^2+2\sum_{i=1}^3|\nabla\phi_i|^2
=-2\lambda_1(\Sigma)+2\sum_{i=1}^3|\nabla\phi_i|^2.
\end{aligned}
\label{eq:7.42}
\end{equation}
因此
\begin{equation}
\lambda_1(\Sigma)=2\sum_{i=1}^3|\nabla\phi_i|^2.
\label{eq:7.43}
\end{equation}
现在, 由于 $\Phi$ 共形, 由 \eqref{eq:7.43} 容易看出
\begin{equation}
2(d\Phi)^{\star}g_{\bS^2}=\lambda_1(\Sigma)g_{\Sigma},
\label{eq:7.44}
\end{equation}
这就证明了所要的结论.
\end{proof}

\begin{theorem}[Yang--Yau, \cite{YgYa}]
\label{thm:7.7}
设 $\Sigma$ 是带 Riemann 度量的曲面, $\Phi:\Sigma\to\bS^2\subset\R^3$ 是
满足 $\deg(\Phi)>0$ 的 (分支) 共形映射, 则
\begin{equation}
\lambda_1(\Sigma)\Area(\Sigma)
\leq\deg(\Phi)\lambda_1(\bS^2)\Area(\bS^2)
=8\pi\deg(\Phi),
\label{eq:7.45}
\end{equation}
等号成立当且仅当 $\Sigma$ 等距于圆球面.
\end{theorem}

\begin{proof}
对 $i=1,2,3$ 置 $\phi_i=x_i\circ\Phi$. 由 Hersch 定理证明开头的论证, 且
因 $\Phi$ 非常值, 可设 \eqref{eq:7.36} 对每个 $i$ 成立. 现在与 Hersch
定理证明的唯一差别在于: 代替 \eqref{eq:7.39}, 我们有
\begin{equation}
\int_{\Sigma}|\nabla\phi_i|^2
=\deg(\Phi)\int_{\bS^2}|\nabla x_i|^2
=\deg(\Phi)\lambda_1(\bS^2)\int_{\bS^2}x_i^2.
\label{eq:7.46}
\end{equation}
完成这一改动即证得定理.
\end{proof}

\begin{lemma}[参看 \cite{FK} 第 124 页]
\label{lem:7.8}
设 $\Sigma$ 是亏格为 $g$ 的闭曲面. 则存在 (分支) 共形映射
$\Phi:\Sigma\to\bS^2$, 其度至多
\begin{equation}
\rho(g)\leq
\begin{cases}
g/2+1, & \text{若 } g \text{ 为偶数},\\
g/2+3/2, & \text{若 } g \text{ 为奇数}.
\end{cases}
\label{eq:7.47}
\end{equation}
\end{lemma}

把定理~\ref{thm:7.7} 与引理~\ref{lem:7.8} 结合起来, 得到:

\begin{theorem}[Yang--Yau, \cite{YgYa}]
\label{thm:7.9}
设 $\Sigma^2$ 是亏格为 $g$ 且带 Riemann 度量的曲面, 则
\begin{equation}
\lambda_1(\Sigma)\Area(\Sigma)
\leq4\pi\bigl(g+5/2-(-1)^g/2\bigr)
\leq8\pi(1+g),
\label{eq:7.48}
\end{equation}
其中第一个不等式取等号当且仅当 $\Sigma$ 是圆球面.
\end{theorem}

已知: 若 $\Sigma$ 是常负曲率的 Riemann 曲面, 则
\begin{equation}
\lambda_1(\Sigma)\Area(\Sigma)\leq\pi g-\pi.
\label{eq:7.49}
\end{equation}

% ================================================================
%  7.3  Reilly 公式
% ================================================================
\section{Reilly 公式}
\label{sec:7.3}

在 \cite{Ry} 中, R. Reilly 发现了 Bochner 公式对于带边界 $\partial\Omega$ 的
区域 $\Omega$ 的一个有用的积分形式.

\begin{lemma}[Reilly, \cite{Ry}]
\label{lem:7.10}
若 $u$ 是有界区域 $\Omega$ 上的光滑函数, 则
\begin{equation}
\begin{aligned}
&\int_{\Omega}\Bigl(\bigl|\nabla^2u\bigr|^2
+\operatorname{Ric}(\nabla u,\nabla u)-(\Delta u)^2\Bigr)\\
&\qquad=\int_{\partial\Omega}\Bigl(A\bigl((\nabla u)^T,(\nabla u)^T\bigr)
-2u_n\Delta_{\partial\Omega}u+Hu_n^2\Bigr),
\end{aligned}
\label{eq:7.50}
\end{equation}
其中 $u_n$ 是法向导数, 而 $H=-A(e_i,e_i)$ 是 $\partial\Omega$ 的平均曲率.
\end{lemma}

\begin{proof}
出发点是 Bochner 公式 (命题~\ref{prop:1.47})
\begin{equation}
\frac12\Delta|\nabla u|^2
=\bigl|\nabla^2u\bigr|^2+\operatorname{Ric}(\nabla u,\nabla u)
+\langle\nabla u,\nabla\Delta u\rangle.
\label{eq:7.51}
\end{equation}
把它与
\begin{equation}
\dive(\Delta u\nabla u)=(\Delta u)^2+\langle\nabla u,\nabla\Delta u\rangle
\label{eq:7.52}
\end{equation}
结合, 得到
\begin{equation}
\dive\Bigl(\frac12\nabla|\nabla u|^2-\Delta u\nabla u\Bigr)
=\bigl|\nabla^2u\bigr|^2+\operatorname{Ric}(\nabla u,\nabla u)-(\Delta u)^2.
\label{eq:7.53}
\end{equation}
因此可应用 Stokes 定理, 得到
\begin{equation}
\begin{aligned}
&\int_{\partial\Omega}
\Bigl\langle\Bigl(\frac12\nabla|\nabla u|^2-\Delta u\nabla u\Bigr),
\partial_n\Bigr\rangle\\
&\qquad=\int_{\Omega}\Bigl(\bigl|\nabla^2u\bigr|^2
+\operatorname{Ric}(\nabla u,\nabla u)-(\Delta u)^2\Bigr),
\end{aligned}
\label{eq:7.54}
\end{equation}
其中 $\partial_n$ 是 $\partial\Omega$ 的向外的单位法向量.

取 $e_i$ 为 $\partial\Omega$ 的正交标架. 利用 Hessian 的对称性, 有
\begin{equation}
\frac12\bigl\langle\nabla|\nabla u|^2,\partial_n\bigr\rangle
=\bigl\langle\nabla_{\partial_n}\nabla u,\nabla u\bigr\rangle
=\bigl\langle\nabla_{\nabla u}\nabla u,\partial_n\bigr\rangle,
\label{eq:7.55}
\end{equation}
于是算得
\begin{equation}
\begin{aligned}
\langle\nabla_{e_i}\nabla u,\partial_n\rangle
&=\bigl\langle\nabla_{e_i}(\nabla u)^T,\partial_n\bigr\rangle
+\bigl\langle\nabla_{e_i}(\nabla u)^\perp,\partial_n\bigr\rangle\\
&=A\bigl(e_i,(\nabla u)^T\bigr)+e_iu_n,
\end{aligned}
\label{eq:7.56}
\end{equation}
其中最后一个等号还用到了
$\langle\nabla_{e_i}\partial_n,\partial_n\rangle=0$. 由此推出
\begin{equation}
\begin{aligned}
\bigl\langle\nabla_{\nabla u}\nabla u,\partial_n\bigr\rangle
&=u_n\bigl\langle\nabla_{\partial_n}\nabla u,\partial_n\bigr\rangle
+A\bigl((\nabla u)^T,(\nabla u)^T\bigr)\\
&\quad+\bigl\langle(\nabla u)^T,(\nabla u_n)^T\bigr\rangle.
\end{aligned}
\label{eq:7.57}
\end{equation}
同样地, 算得
\begin{equation}
\begin{aligned}
\Delta u&=\bigl\langle\nabla_{\partial_n}\nabla u,\partial_n\bigr\rangle
+\bigl\langle\nabla_{e_i}(\nabla u)^T,e_i\bigr\rangle
+\bigl\langle\nabla_{e_i}(\nabla u)^\perp,e_i\bigr\rangle\\
&=\bigl\langle\nabla_{\partial_n}\nabla u,\partial_n\bigr\rangle
+\Delta_{\partial\Omega}u-A(e_i,e_i)u_n.
\end{aligned}
\label{eq:7.58}
\end{equation}
最后, 把以上各式综合起来, 即得边界项的一个表达式:
\begin{equation}
\begin{aligned}
\Bigl(\frac12\partial_n|\nabla u|^2-\Delta uu_n\Bigr)
&=A\bigl((\nabla u)^T,(\nabla u)^T\bigr)
+\bigl\langle(\nabla u)^T,(\nabla u_n)^T\bigr\rangle\\
&\quad-u_n\Delta_{\partial\Omega}u+Hu_n^2,
\end{aligned}
\label{eq:7.59}
\end{equation}
其中 $H=-A(e_i,e_i)$ 是 $\partial\Omega$ 的平均曲率. 再对 $\partial\Omega$
上的项 $\langle(\nabla u)^T,(\nabla u_n)^T\rangle$ 分部积分, 即得引理.
\end{proof}

% ================================================================
%  7.4  Choi--Wang 关于 $\lambda_1$ 的下界
% ================================================================
\section{Choi--Wang 关于 $\lambda_1$ 的下界}
\label{sec:7.4}

本节遵循 Choi 与 Wang 的做法, 用 Reilly 公式证明正 Ricci 曲率流形中闭
嵌入极小超曲面的第一特征值的一个下界.

现在设 $\partial\Omega$ 极小 (即 $H=0$). 令 $\lambda$ 为 $\Delta_{\partial\Omega}$
的第一 Dirichlet 特征值, $f$ 为相应的特征函数, 于是
\begin{equation}
\Delta_{\partial\Omega}f=-\lambda f
\quad\text{且}\quad
\int_{\partial\Omega}f^2=1.
\label{eq:7.60}
\end{equation}
令 $u$ 为 $\Omega$ 上以 $f$ 为边界数据的 Dirichlet 问题的解, 于是
\begin{equation}
\Delta_{\Omega}u=0
\quad\text{且}\quad
u\big|_{\partial\Omega}=f.
\label{eq:7.61}
\end{equation}
我们将在引理~\ref{lem:7.10} 中使用这个函数 $u$. 其中的体积分变成
\begin{equation}
\int_{\Omega}\Bigl(\bigl|\nabla^2u\bigr|^2
+\operatorname{Ric}(\nabla u,\nabla u)\Bigr)
\geq\int_{\Omega}\operatorname{Ric}(\nabla u,\nabla u).
\label{eq:7.62}
\end{equation}
与此同时, 边界积分变成
\begin{equation}
\begin{aligned}
&\int_{\partial\Omega}\Bigl(A\bigl((\nabla u)^T,(\nabla u)^T\bigr)
-2u_n\Delta_{\partial\Omega}u\Bigr)
=\int_{\partial\Omega}\Bigl(A\bigl((\nabla u)^T,(\nabla u)^T\bigr)
+2\lambda u_nu\Bigr)\\
&\qquad=\int_{\partial\Omega}A\bigl((\nabla u)^T,(\nabla u)^T\bigr)
+2\lambda\int_{\Omega}|\nabla u|^2.
\end{aligned}
\label{eq:7.63}
\end{equation}
把边界项与内部项结合起来, 给出
\begin{equation}
2\lambda\int_{\Omega}|\nabla u|^2
\geq(\min\operatorname{Ric})\int_{\Omega}|\nabla u|^2
-\int_{\partial\Omega}A\bigl((\nabla u)^T,(\nabla u)^T\bigr).
\label{eq:7.64}
\end{equation}
为完成论证, 注意若把 $\Omega$ 换为其补集, 则 $A$ 这一项变号. 于是得到
Choi 与 Wang 的定理 \cite{CiWa}:

\begin{theorem}[Choi--Wang, \cite{CiWa}]
\label{thm:7.11}
若 $M$ 是满足 $\operatorname{Ric}_M>k$ 的单连通闭流形,
$\Sigma\subset M$ 是闭嵌入极小超曲面, 则 $\Sigma$ 的第一特征值至少为
$\frac k2$.
\end{theorem}

最后我们指出, Choi 与 Schoen \cite{CiSc} 用覆叠论证把定理~\ref{thm:7.11}
推广到了一切具有正 Ricci 曲率的闭流形 $M$ (即他们去掉了 $M$ 单连通的
假设).

% ================================================================
%  7.5  带先验界的紧性定理
% ================================================================
\section{带先验界的紧性定理}
\label{sec:7.5}

本节将证明 Choi 与 Schoen \cite{CiSc} 的下述紧性定理:

\begin{theorem}[Choi--Schoen, \cite{CiSc}]
\label{thm:7.12}
正 Ricci 曲率的光滑闭三流形 $M^3$ 中, 亏格为 $g$ 的闭嵌入极小曲面所成
之空间在光滑拓扑下是紧的.
\end{theorem}

这一结果的证明利用 Ricci 曲率的正性, 得到嵌入极小曲面面积以其亏格表示
的先验上界. 具体地说, 把定理~\ref{thm:7.11} 与前面的定理~\ref{thm:7.9}
结合, Choi 与 Wang 得到下述推论:

\begin{corollary}[Choi--Wang, \cite{CiWa}]
\label{cor:7.13}
若 $M^3$ 满足 $\operatorname{Ric}_M\geq\Lambda>0$, 且
$\Sigma_g^2\subset M$ 是亏格为 $g$ 的闭嵌入极小曲面, 则
\begin{equation}
\Area(\Sigma_g)\leq\frac{16\pi(1+g)}{\Lambda}.
\label{eq:7.65}
\end{equation}
\end{corollary}

对 \eqref{eq:7.7} 积分并应用 Gauss--Bonnet 公式, 得
\begin{equation}
\int_{\Sigma}|A|^2
\leq2k\Area(\Sigma_g)+8\pi(g-1)
\leq\frac{32\pi k(1+g)}{\Lambda}+8\pi(g-1),
\label{eq:7.66}
\end{equation}
其中第二个不等式由 \eqref{eq:7.65} 得到. \eqref{eq:7.65} 与 \eqref{eq:7.66}
的意义在于: $\Sigma$ 的面积与全曲率都以亏格为界一致有界. 这对于得到
光滑紧性定理是本质的. 注意单是面积界就足以给出整变分流形空间中的紧性
(参看\secref{sec:7.1}).

下一个命题说明如何把界 \eqref{eq:7.65} 与 \eqref{eq:7.66} 同定理~\ref{thm:2.2}
结合起来, 得到极小曲面的紧性定理 (参看 \cite{CiSc}; 并参看 M. T. Anderson
\cite{An} 与 B. White \cite{Wh1}).

\begin{proposition}
\label{prop:7.14}
设 $M^3$ 是闭三流形, $\Sigma_i\subset M$ 是一列亏格为 $g$ 的闭嵌入极小
曲面, 满足
\begin{equation}
\Area(\Sigma_i)\leq C_1
\label{eq:7.67}
\end{equation}
与
\begin{equation}
\int_{\Sigma_i}\bigl|A_{\Sigma_i}\bigr|^2\leq C_2.
\label{eq:7.68}
\end{equation}
则存在有限点集 $\mathcal{S}\subset M$ 与子列 $\Sigma_{i'}$, 使得
$\Sigma_{i'}$ 在 $M\setminus\mathcal{S}$ 的紧子集上按 $C^\ell$ 拓扑 (任意
$\ell<\infty$) 一致收敛到极小曲面 $\Sigma\subset M$. 该子列还按 (外蕴)
Hausdorff 距离收敛到 $\Sigma$. 而 $\Sigma$ 在 $M$ 中光滑、亏格至多 $g$,
且满足 \eqref{eq:7.67} 与 \eqref{eq:7.68}.
\end{proposition}

\begin{proof}
我们只对 $\ell=2$ 给出证明, 因为一般情形可由标准椭圆理论推出. 本证明
中, $\epsilon=\epsilon(M)>0$ 与 $r_0=r_0(M)>0$ 都取自定理~\ref{thm:2.2}. 取 $r_1$
使得由 \eqref{eq:7.5}, 对 $s<t\leq r_1$ 及任意极小曲面 $\Gamma\subset M$
有
\begin{equation}
\frac12s^{-2}\Area(B_s\cap\Gamma)
\leq t^{-2}\Area(B_t\cap\Sigma).
\label{eq:7.69}
\end{equation}

为找出集合 $\mathcal{S}$, 定义测度 $\nu_i$ 为
\begin{equation}
\nu_i(U)=\int_{U\cap\Sigma_i}|A_i|^2\leq C_2,
\label{eq:7.70}
\end{equation}
其中 $U\subset M$, $A_i=A_{\Sigma_i}$. Radon 测度的一般紧性定理 (即定理~\ref{thm:3.2})
蕴含存在子列 $\nu_{\beta_i}$ 弱收敛到 Radon 测度 $\nu$, 且
\begin{equation}
\nu(M)\leq C_2.
\label{eq:7.71}
\end{equation}
为记号简便, 把 $\nu_i$ 换成 $\nu_{\beta_i}$. 定义集合
\[
\mathcal{S}=\{x\in M\mid\nu(x)\geq\epsilon\}.
\]
由 \eqref{eq:7.71} 立即得知 $\mathcal{S}$ 至多含 $C_2\epsilon^{-1}$ 个点.

给定任意 $y\in M\setminus\mathcal{S}$, 有 $\nu(y)<\epsilon$. 由于 $\nu$ 是
Radon 测度, 从而是 Borel 正则的, 故存在某个 $0<10s<\min\{r_0,r_1\}$
(依赖于 $y$), 使得
\begin{equation}
\nu(B_{10s}(y))<\epsilon.
\label{eq:7.72}
\end{equation}
由于 $\nu_i\to\nu$, \eqref{eq:7.72} 蕴含对充分大的 $i$
\begin{equation}
\int_{B_{10s}(y)\cap\Sigma_i}|A_i|^2<\epsilon.
\label{eq:7.73}
\end{equation}
这使我们可以对每个 $B_{10s}(y)\cap\Sigma_i$ 一致地应用定理~\ref{thm:2.2}. 由此可知,
对充分大的 $i$ 及 $z\in B_{5s}(y)\cap\Sigma_i$,
\begin{equation}
25s^2|A_i|^2(z)\leq1.
\label{eq:7.74}
\end{equation}
由 (引理~\ref{lem:2.4} 的一个小变形), \eqref{eq:7.74} 蕴含: 对每个
$z\in B_s(y)\cap\Sigma_i$, $B_s(z)\cap\Sigma_i$ 中含 $z$ 的连通分支是
$T_z\Sigma_i$ 中某个区域 $U_z^i$ 上函数 $u_i^z$ 的图, 且
$|\nabla u_i^z|\leq1$, $2s|\nabla^2u_i^z|\leq1$. $|\nabla u_i^z|$ 与
$|\nabla^2u_i^z|$ 的这些界蕴含 $u_i^z$ 满足一个具有 Lipschitz 系数的一致
椭圆方程 (参看 \eqref{eq:7.19}). 通常的椭圆估计 (例如参看 \cite{GiTr}
推论 6.3) 给出某个 $\alpha>0$ 以及 $u_i^z$ 在 $B_{s/2}\subset T_z\Sigma_i$
上一致的 $C^{2,\alpha}$ 估计.

由单调性 (即 \eqref{eq:7.69}), $B_s(y)\cap\Sigma_i$ 的每个与
$B_{s/2}(y)$ 相交的连通分支面积至少为 $\frac\pi8s^2$. 记这样的分支个数为
$c_y$. 由单调性、\eqref{eq:7.69} 与 \eqref{eq:7.67}, 得
\begin{equation}
\frac\pi{16}r_1^2c_y\leq C_1.
\label{eq:7.75}
\end{equation}
特别地, $B_s(y)\cap\Sigma_i$ 中与 $B_{s/2}(y)$ 相交的连通分支个数不依赖
于 $i$ 与 $y$ 而有界.

既然对分支个数有一致估计, 又对各个图有一致的 $C^{2,\alpha}$ 估计,
Arzelà--Ascoli 定理给出另一子列 $\eta_i$, 它在 $B_s(y)$ 中按
$C^{2,\frac\alpha2}$ 拓扑一致收敛. 由于可用可数多个这样的球覆盖
$M\setminus\mathcal{S}$, 用对角线论证即可完成到某个 (可能为浸入的) 曲面
$\Sigma$ 的收敛, 该曲面在 $M\setminus\mathcal{S}$ 中光滑. 这也蕴含
$\Sigma$ 满足 \eqref{eq:7.67} 与 \eqref{eq:7.68}. 此外, 一致的
$C^{2,\frac\alpha2}$ 估计蕴含 $\Sigma$ 满足同一微分方程, 即 $\Sigma$ 是
极小的.

我们还可设 $\Sigma_i$ 作为整变分流形收敛到支撑在 $\Sigma$ 上的变分流形.
由常值性定理, 它必是 $\Sigma$ 的某个倍数 (\cite{Si4} 定理 41.1). 这一
收敛蕴含单调性 (即 \eqref{eq:7.69}) 适用于极限 $\Sigma$. 把它与面积界
结合, 对任意 $y\in M$ 与 $r<r_1$ 有
\begin{equation}
r^{-2}\Area(B_r(y)\cap\Sigma)\leq2C_1.
\label{eq:7.76}
\end{equation}
单调性还蕴含 $\Sigma_i$ 按 Hausdorff 距离收敛到 $\Sigma$. 为看出这一点,
注意若 $y\in\Sigma_i$ 满足 $\dist(y,\Sigma)>2\delta$, 则单调性蕴含
$B_\delta(y)\cap\Sigma_i$ 的面积至少为 $C'\delta^2$. 由于变分流形收敛蕴含
面积测度收敛, 可见 $\Sigma_i$ 必按 Hausdorff 距离收敛到 $\Sigma$.

由于每个 $\Sigma_i$ 都是嵌入的, 而收敛是光滑的, $\Sigma$ 不会自交. 于是
定理~\ref{thm:7.3} 的局部描述蕴含 $\Sigma$ 必是嵌入的.

剩下要证 $\Sigma\cup\mathcal{S}$ 是嵌入极小曲面. 换言之, 必须证明每个
$x\in\mathcal{S}$ 都是 $\Sigma$ 的可去奇点. 记 $A$ 为 $\Sigma$ 的第二基本
形式, 于是 (利用 \eqref{eq:7.68}) $|A|^2$ 是 $\Sigma$ 上的 $L^1$ 函数.
单调收敛定理蕴含
\begin{equation}
\lim_{r\to0}\int_{B_r(x)\cap\Sigma}|A|^2=0.
\label{eq:7.77}
\end{equation}
因此, 给定任意 $0<\delta<1$, 存在某个 $0<r_x<r_0$, 使得
\begin{equation}
\int_{B_{2r_x}(x)\cap\Sigma}|A|^2<\delta\epsilon.
\label{eq:7.78}
\end{equation}
把定理~\ref{thm:2.2} 用于 $\Sigma$ 自身, 若 $r<r_x$ 且
$z\in B_r(x)\setminus B_{r/2}(x)$, 则
\begin{equation}
r^2|A|^2(z)\leq4\delta.
\label{eq:7.79}
\end{equation}
暂时固定 $r<r_x$ 与 $z_1\in\partial B_{3r/4}(x)$. 引理~\ref{lem:2.4} 的一个小变形
与 \eqref{eq:2.54} 蕴含: 当 $\delta$ 充分小时, $B_{r/4}(z_1)\cap\Sigma$ 中含 $z$
的分支是 $T_z\Sigma$ 上一个梯度以 $C\sqrt\delta$ 为界的极小图. 若对该极小
图中的某个
\[
z_2\in\partial B_{3r/4}(x)\cap\partial B_{r/8}(z_1)
\]
重复这一论证, 则可知
\[
B_{r/4}(z_2)\cap\Sigma
\]
中含 $z$ 的连通分支也是梯度有界的极小图. 面积界 \eqref{eq:7.76} 蕴含:
沿 $\partial B_{3r/4}(x)$ 把这个论证迭代大约 $C_1$ 次之后必定闭合. 取
$\delta$ 充分小 (只依赖于 $C_1$), 可知
$(B_r(x)\setminus B_{r/2}(x))\cap\Sigma$ 中含 $z_1$ 的连通分支是某个固定
切平面上的图, 其梯度以 $C\sqrt\delta$ 为界, Hessian 以 $2\sqrt\delta r^{-1}$
为界 (参看引理~\ref{lem:2.4}).

取如上足够小的 $\delta>0$ 与 $r_x>0$, 令 $\Sigma_x$ 为
$B_{r_x}(x)\cap\Sigma$ 的一个分支, 且 $x\in\bar\Sigma_x$. 上述讨论表明
$\Sigma_x$ 是某个固定切平面上函数 $u$ 的极小图, 且 $|\nabla u|\leq C\sqrt\delta$.
下面来说明 $\nabla u$ 在 $x$ 处有极限. 为此注意: 对任意 $\delta_c>0$, 可
仿照 \eqref{eq:7.78} 与 \eqref{eq:7.79} 论证, 找到某个 $0<r_c<r_x$,
使得对任意 $r<r_c$ 与 $z\in\partial B_r(x)\cap\Sigma_x$ 有
\begin{equation}
r^2\bigl|\Hess_u\bigr|\leq4r^2|A|^2(z)\leq16\delta_c.
\label{eq:7.80}
\end{equation}
沿 $\partial B_r(x)$ 积分, 并利用 $\partial B_r(x)\cap\Sigma_x$ 是梯度
有界的图这一事实, 可知
\begin{equation}
\sup_{z_1,z_2\in\partial B_r(x)\cap\Sigma_x}
|\nabla u(z_1)-\nabla u(z_2)|\leq16\pi\sqrt{\delta_c}.
\label{eq:7.81}
\end{equation}
由此立即推出 $\nabla u$ 在 $x$ 处有极限. 特别地, $u$ 可延拓为极小曲面
方程的 $C^1$ 解, 且梯度一致地小. 现在可应用引理~\ref{lem:7.2} 断定
$\Sigma_x\cup\{x\}$ 是光滑极小曲面.

对 $\mathcal{S}$ 的每个元素重复这一做法, 可见 $\Sigma\cup\mathcal{S}$ 是
光滑极小曲面, 且在 $\mathcal{S}$ 之外是嵌入的. 而局部描述 (即定理~\ref{thm:7.3})
蕴含它处处嵌入. 这就完成了证明.
\end{proof}

现在可以证明定理~\ref{thm:7.12} 了. 由于在这一情形下有面积界与全曲率界,
命题~\ref{prop:7.14} 蕴含: 任意一列亏格固定的闭嵌入极小曲面都有子列,
它在有限多个点之外收敛到光滑嵌入极小曲面. 剩下主要是要证明, 越过这些点
收敛依然成立. 这有多种做法, 我们采用 \cite{CiSc} 原始证明中关于
$\lambda_1$ 下界的论证.

\begin{proof}[定理~\ref{thm:7.12} 的证明]
由 Myers 定理, 由于 $\operatorname{Ric}_M>0$, $M$ 的 $\pi_1$ 有限; 于是
取有限覆叠后可设 $M$ 单连通. 只需证明: 给定 $M$ 中任意一列亏格为 $g$ 的
闭嵌入极小曲面 $\Sigma_i$, 存在子列按 $C^{2,\alpha}$ 拓扑 (某个 $\alpha>0$)
收敛.

由于 $\operatorname{Ric}_M>0$, \eqref{eq:7.65} 与 \eqref{eq:7.66} 成立.
因此可应用命题~\ref{prop:7.14}. 于是存在有限点集
$\mathcal{S}\subset M$ 与子列 $\Sigma_{i'}$, 使它在 $M\setminus\mathcal{S}$
的紧子集上按 $C^{2,\alpha}$ 拓扑一致收敛到极小曲面 $\Sigma\subset M$.
$\Sigma$ 在 $M$ 中光滑且嵌入, 亏格至多 $g$, 并满足 \eqref{eq:7.65} 与
\eqref{eq:7.66}. 由于 $M$ 单连通而 $\Sigma$ 嵌入, $\Sigma$ 的法丛是平凡的.

剩下要证收敛越过点集 $\mathcal{S}$ 仍是光滑的. 我们将看到, 这等价于
证明收敛是单重的 (这一等价也可由 Allard 正则性定理推出). 先用定理
\ref{thm:7.11} 来说明收敛是单重的.

给定任意 $\epsilon>0$, 存在 $i_0=i_0(\epsilon)$, 使得对任意 $i>i_0$ 有
\begin{equation}
\Sigma_i\setminus\bigcup_{x_j\in\mathcal{S}}B_{\epsilon^2}(x_j)
=\bigcup_{n=1}^m\Sigma_{i,n},
\label{eq:7.82}
\end{equation}
其中每个 $\Sigma_{i,n}$ 都是 $\Sigma$ 上梯度有界的不相交极小图.
$\Sigma_i$ 嵌入以及 $\Sigma$ 的法丛平凡这两点蕴含我们可以取``最上面''的
一个图, 记之为 $\Sigma_{i,1}$.

于是, 若收敛不是单重的, 则当 $i$ 大时极小曲面 $\Sigma_i$ 由若干大片图
在小集合上连接而成. 利用这一描述, 我们将构造一个除这些小集合附近
之外为常值的试验函数. 下面将证明它违背 $\lambda_1$ 的下界.

在 $\Sigma_i$ 上定义 Lipschitz 函数 $\phi$ 为
\[
\begin{aligned}
\phi(x)&=1 &&\text{对 } x\in\Sigma_{i,1}\setminus
  \bigcup_{x_j\in\mathcal{S}}B_\epsilon(x_j),\\
\phi(x)&=\frac{\log\dist_M(x,x_j)-2\log\epsilon}{\log\epsilon-2\log\epsilon}
 &&\text{对 } x\in\Sigma_{i,1}\cap B_\epsilon(x_j)\setminus
  B_{\epsilon^2}(x_j),\\
\phi(x)&=0 &&\text{对 } x\in\Sigma_{i,1}\cap
  \Bigl(\bigcup_{x_j\in\mathcal{S}}B_{\epsilon^2}(x_j)\Bigr),\\
\phi(x)&=-\frac{\log\dist_M(x,x_j)-2\log\epsilon}{\log\epsilon-2\log\epsilon}
 &&\text{对 } n>1,\ x\in\Sigma_{i,n}\cap B_\epsilon(x_j)\setminus
  B_{\epsilon^2}(x_j),\\
\phi(x)&=-1 &&\text{对 } n>1 \text{ 及 } x\in\Sigma_{i,n}\setminus
  \bigcup_{x_j\in\mathcal{S}}B_\epsilon(x_j).
\end{aligned}
\]
置 $c=\int_{\Sigma_i}\phi$, 于是 $\psi=\phi-c$ 的积分为零. 显然, 至少在
某一层上有
\begin{equation}
|\psi|\geq1+|c|\geq1
\label{eq:7.83}
\end{equation}
从而
\begin{equation}
\lim_{\epsilon\to0}\int_{\Sigma_i}\psi^2\geq\Area(\Sigma).
\label{eq:7.84}
\end{equation}
另一方面, $\nabla\psi=\nabla\phi$, 故由余面积公式 (即 \eqref{eq:1.59}) 得
\begin{equation}
\begin{aligned}
\int_{\Sigma_i}|\nabla\psi|^2
&\leq\frac1{(\log\epsilon)^2}\sum_{x_j\in\mathcal{S}}
\int_{B_\epsilon(x_j)\setminus B_{\epsilon^2}(x_j)}
\frac{|\nabla r_j|^2}{r_j^2}\\
&=\frac1{(\log\epsilon)^2}\sum_{x_j\in\mathcal{S}}
\int_{t=\epsilon^2}^{\epsilon}
\int_{\{r_j=t\}\cap\Sigma_i}\frac{|\nabla r_j|}{t^2},
\end{aligned}
\label{eq:7.85}
\end{equation}
其中 $r_j(x)=\dist(x,x_j)$. 用 Stokes 定理与单调性可把它估计为
\begin{equation}
\begin{aligned}
\int_{\{r_j=t\}\cap\Sigma_i}\frac{|\nabla r_j|}{t^2}
&=\frac12t^{-3}\int_{B_t(x_j)\cap\Sigma_i}\Delta r_j^2\\
&\leq3t^{-2}\Area(B_t(x_j)\cap\Sigma_i)\leq3C_1t^{-1}.
\end{aligned}
\label{eq:7.86}
\end{equation}
把 \eqref{eq:7.86} 代入 \eqref{eq:7.85} 并积分,
\begin{equation}
\int_{\Sigma_i}|\nabla\psi|^2
\leq3C_1|\mathcal{S}|\frac1{(\log\epsilon)^2}
\int_{\epsilon^2}^{\epsilon}\frac{dt}t
=3C_1|\mathcal{S}|\frac1{\log\epsilon}.
\label{eq:7.87}
\end{equation}
因此, 取 $\epsilon$ 充分小, 试验函数 $\psi$ 就违背定理~\ref{thm:7.11}.
这一矛盾表明 $m=1$, 即收敛是单重的.

利用这一点, 我们将在每个 $x_j$ 的邻域内得到 $\Sigma_i$ 的一致曲率估计.
如前, 把它与 Arzelà--Ascoli 定理结合, 将蕴含即使在 $\mathcal{S}$ 上也是
光滑收敛.

给定 $x_j\in\mathcal{S}$ 与 $0<\epsilon$, 由于 $\Sigma$ 光滑, 可取 $0<s$
使 $B_{2s}(x_j)\cap\Sigma$ 是梯度以 $\epsilon$ 为界的极小图. 对充分大的
$i$ 可设
\[
(B_{2s}(x_j)\setminus B_s(x_j))\cap\Sigma_i
\]
是 $\Sigma$ 上函数 $u_i$ 的图, 且
\[
s^{-1}|u|+|\nabla u|\leq2\epsilon.
\]
由于 $\Sigma_i\to\Sigma$ 按 Hausdorff 距离, 对充分大的 $i$ 还有 $\Sigma_i$
含于 $\frac s4$ 管状邻域之中. 注意: 正是收敛的单重性被用来断定这是单个
图. 由于 $\Sigma_i$ 含于该管状邻域, $B_{2s}(x_j)\cap\Sigma_i$ 的边界就是
$\partial B_{2s}(x_j)\cap\Sigma_i$. 利用这一梯度小的图描述, 并对
$\Delta_{\Sigma_i}r_j^2$ 应用 Stokes 定理, 得
\begin{equation}
\frac{\Area(B_{2s}(x_j)\cap\Sigma_i)}{4\pi s^2}\leq1+C_3\epsilon.
\label{eq:7.88}
\end{equation}
由单调性, \eqref{eq:7.88} 蕴含对任意 $z\in B_s(x_j)\cap\Sigma_i$ 有
\begin{equation}
\frac{\Area(B_s(z)\cap\Sigma_i)}{\pi s^2}\leq1+C_4\epsilon.
\label{eq:7.89}
\end{equation}
最后, Allard 正则性定理的光滑版本 (即定理~\ref{thm:2.15}) 给出当 $\epsilon>0$
充分小时一致的曲率估计, 证毕.
\end{proof}

% ================================================================
%  7.6  正质量定理
% ================================================================
\section{正质量定理}
\label{sec:7.6}

本节将证明 Schoen 与 Yau 的正质量定理 \cite{ScYa1}. 该证明大致遵循
\cite{ScYa1}, 是极小曲面理论在广义相对论中的一个应用.

设 $M$ 是渐近平坦的, 且 $\operatorname{Scal}_{M^3}\geq0$. 我们假定 $M$
只有一个端; 一般情形的改动用同样的办法即可. 具体地说, 存在紧集
$\Omega\subset M$, 使 $M\setminus\Omega$ 即为 $\R^3\setminus B_R(0)$,
且 $M\setminus\Omega$ 上的度量可写为
\begin{equation}
g_{ij}=\Bigl(1+\frac{\mathcal{M}}{2|x|}\Bigr)^4\delta_{ij}+p_{ij},
\label{eq:7.90}
\end{equation}
其中
\begin{equation}
|x|^2|p_{ij}|+|x|^3|Dp_{ij}|+|x|^4|D^2p_{ij}|\leq C.
\label{eq:7.91}
\end{equation}
$\mathcal{M}$ 是 $M$ 的\emph{质量}. 当 $p_{ij}\equiv0$ 时,
$M\setminus\Omega$ 等距于质量为 $\mathcal{M}$ 的 Schwarzschild 时空中的
一个常时 (类空) 切片.

称张量 $h$ 为 $O(|x|^{-p})$ 的, 若
\[
|x|^p|h|+|x|^{p+1}|Dh|\leq C.
\]
例如,
\begin{equation}
\begin{aligned}
g_{ij}&=\bigl(1+2\mathcal{M}/|x|\bigr)\delta_{ij}+O\bigl(|x|^{-2}\bigr),\\
\sqrt{g}&\equiv\sqrt{\det g_{ij}}
=1+3\mathcal{M}|x|^{-1}+O\bigl(|x|^{-2}\bigr).
\end{aligned}
\label{eq:7.92}
\end{equation}

现在可以叙述 Schoen 与 Yau 的正质量定理了:

\begin{theorem}[正质量定理, \cite{ScYa1}]
\label{thm:7.15}
若 $M$ 如上所述, 则 $\mathcal{M}\geq0$.
\end{theorem}

\begin{theorem}[唯一性, \cite{ScYa1}]
\label{thm:7.16}
若 $|\nabla^3p_{ij}|=O(|x|^{-5})$, 且 \ref{thm:7.15} 中
$\mathcal{M}=0$, 则 $M=\R^3$.
\end{theorem}

我们只叙述了三维的情形, 但 Schoen 与 Yau 证明了直到七维的类似结果; 维数
的限制来自极小化超曲面的正则性理论, 该理论容许余维七的奇点. Witten
\cite{Wi} 后来用旋量给出了 $\mathcal{M}\geq0$ 的另一个证明, 它对一切维数
都成立, 但只在 $M$ 具有自旋结构 ($M$ 上的一个拓扑条件) 时才可用.

最后我们指出, Penrose 猜测: 当 $M$ 含闭极小曲面时有一个更强的不等式,
即质量 $\mathcal{M}$ 以 $\sqrt{\frac{\mathcal{A}}{16\pi}}$ 为下界, 其中
$\mathcal{A}$ 是``最外层''闭极小曲面的面积. 这一 Penrose 不等式由
T. Ilmanen 与 G. Huiskin \cite{HuI} 用逆平均曲率流证明. 不久之后, H. Bray
\cite{Br} 用正质量定理给出了一个十分不同的证明, 并在最外层极小曲面有
多个分支时给出更强的估计.

\mtpara{正质量定理的证明}
我们将通过证明 $\mathcal{M}<0$ 会导致矛盾来证明定理~\ref{thm:7.15}.

先做一些计算. 回忆 Christoffel 符号 $\Gamma_{jk}^i$ 由
\begin{equation}
\nabla_{\partial_i}\partial_j=\Gamma_{ij}^k\partial_k
\label{eq:7.93}
\end{equation}
定义, 于是函数 $f$ 的 Hessian 由
\begin{equation}
\nabla_{ij}f=f_{ij}-\Gamma_{ij}^kf_k
\label{eq:7.94}
\end{equation}
给出. 利用 \eqref{eq:7.90}、\eqref{eq:7.91} 与
$\partial_i|x|=x_i/|x|$, 算得
\begin{equation}
\begin{aligned}
\Gamma_{ij}^k&=(1/2)g^{km}\bigl(g_{im,j}+g_{jm,i}-g_{ij,m}\bigr)\\
&=-\mathcal{M}\bigl(x_j\delta_{ik}+x_i\delta_{jk}-x_k\delta_{ij}\bigr)
|x|^{-3}+O\bigl(|x|^{-3}\bigr).
\end{aligned}
\label{eq:7.95}
\end{equation}
由此可见, 对这种度量有
\begin{equation}
|x|^2\bigl|\Gamma_{jk}^i\bigr|+|x|^3\bigl|R_{jkl}^i\bigr|\leq C.
\label{eq:7.96}
\end{equation}

\begin{lemma}
\label{lem:7.17}
若 $\mathcal{M}<0$, $\operatorname{Scal}_M\geq0$, 且 $M$ 渐近平坦, 则可把
$g$ 换成一个新的渐近平坦度量, 使得 $\operatorname{Scal}_M\geq0$,
$\operatorname{Scal}_M>0$ 于 $|x|$ 大处, 且质量为 $\mathcal{M}/2$.
\end{lemma}

\begin{proof}[证明概要]
容易算得 $\Delta|x|^{-1}<0$ 于 $|x|$ 大处. 这使我们能把函数``磨圆'',
得到 $\phi>0$ 满足 $\Delta\phi\geq0$, $\Delta\phi>0$ 于 $|x|$ 大处, 且
$\phi$ 渐近于 $1-\mathcal{M}/(4|x|)$. 把 $g$ 换成 $\phi^4g$ 即可.
\end{proof}

\begin{lemma}
\label{lem:7.18}
若 $\mathcal{M}<0$ 且 $M$ 渐近平坦, 则存在 $R_0,h>0$, 使得对 $r>R_0$,
集合
\begin{equation}
C_r=\bigl\{|x|^2\leq r^2,\ -h\leq x_3\leq h\bigr\}
\label{eq:7.97}
\end{equation}
的边界是严格平均凸的.
\end{lemma}

\begin{proof}
利用 \eqref{eq:7.94} 与 \eqref{eq:7.95}, 算得
\begin{equation}
\nabla_{ij}|x|^2=2\delta_{ij}-2x_k\Gamma_{ij}^k
=2\delta_{ij}+O\bigl(|x|^{-1}\bigr).
\label{eq:7.98}
\end{equation}
于是 $|x|^2$ 在 $r>R_0$ 时严格凸; 特别地, $\{|x|^2=r^2\}$ 在 $r>R_0$
时严格凸. 设 $\nu$ 是切片 $\{x_3=t\}$ 的单位法向量, 则
\begin{equation}
\nu=\bigl(1-2\mathcal{M}/|x|\bigr)\partial_3+O\bigl(|x|^{-2}\bigr).
\label{eq:7.99}
\end{equation}
由于 $\sqrt{g}=1+3\mathcal{M}/|x|+O(|x|^{-2})$ 且 $|x|_i=x_i/|x|$,
算得
\begin{equation}
\begin{aligned}
\dive(\nu)&=\sqrt{g}^{-1}\bigl(\nu^i\sqrt{g}\bigr)_i
=\bigl(1+O(|x|^{-1})\bigl(\bigl(1-2\mathcal{M}/|x|\bigr)\delta_{i3}
+O(|x|^{-2})\bigr)_i\\
&=-2\mathcal{M}x_3|x|^{-3}+O\bigl(|x|^{-3}\bigr).
\end{aligned}
\label{eq:7.100}
\end{equation}
现在利用 $\mathcal{M}<0$, 可知存在 $h>0$, 使``三明治''
\begin{equation}
\{-h\leq x_3\leq h\}
\label{eq:7.101}
\end{equation}
的边界平均凸.
\end{proof}

\begin{proof}[定理~\ref{thm:7.15} 的证明]
由引理~\ref{lem:7.17}, 可设 $\mathcal{M}<0$,
$\operatorname{Scal}_M\geq0$, 且 $\operatorname{Scal}_M>0$ 于
$|x|>R$. 我们将构造一个完备的渐近平坦稳定极小曲面 $\Sigma$, 它是
Plateau 问题解的一个极限. 取引理~\ref{lem:7.18} 中的 $h,R_0$, 并固定
一条曲线
\begin{equation}
\sigma\subset\{-h\leq x_3\leq h\}\setminus B_R
\label{eq:7.102}
\end{equation}
它从 $\{x_3=-h\}$ 走到 $\{x_3=h\}$. 对 $s>R_0$, 置
\begin{equation}
\gamma_s=\bigl\{x_1^2+x_2^2=s^2,\ x_3=0\bigr\}.
\label{eq:7.103}
\end{equation}
$\gamma_s$ 与 $\sigma$ 的环绕数非零.

解 Plateau 问题 (参看第 4 章与第 6 章), 得到极小化面积的嵌入极小圆盘
\begin{equation}
\Sigma_s\subset C_s,
\label{eq:7.104}
\end{equation}
且 $\partial\Sigma_s=\gamma_s$. 由环绕论证,
\[
\Sigma_s\cap\sigma\neq\emptyset.
\]
由于各 $\Sigma_s$ 极小化面积, 标准的比较论证给出了一致的密度界: 若
$y\in\Omega_s$, $t>R+2h+|y|$ 且 $B_t(y)\subset C_s$, 则
\begin{equation}
\begin{aligned}
&\Area(B_t(y)\cap\Sigma_s)\leq\Area(C_s\cap\partial B_t(y))\\
&\quad+\max\bigl\{\Area(\{x_3=h\}\cap B_t(y)),
\Area(\{x_3=-h\}\cap B_t(y))\bigr\}\\
&\quad\leq4h\pi t+O(t)+\pi t^2+O(t)\leq\pi t^2+O(t).
\end{aligned}
\label{eq:7.105}
\end{equation}

最后, 稳定曲面的曲率估计 (参看第 2 章) 给出
\begin{equation}
|x|^2|A(x)|^2\leq C.
\label{eq:7.106}
\end{equation}
由于 $\Sigma\subset\{|x_3|\leq h\}$, \eqref{eq:7.106} 蕴含 (对 $|x|$
大) $\Sigma$ 的每个连通分支都是梯度很小的图; 下面我们把它加强. 由
\eqref{eq:7.94} 与 \eqref{eq:7.95},
\begin{equation}
\nabla_{ij}x_3=-\Gamma_{ij}^3=(1+x_3)O\bigl(|x|^{-3}\bigr).
\label{eq:7.107}
\end{equation}
由于 $\Sigma$ 极小且 $|x_3|\leq h$ 于 $\Sigma$ 上, \eqref{eq:7.107} 蕴含
\begin{equation}
\Delta_{\Sigma}x_3=O\bigl(|x|^{-3}\bigr).
\label{eq:7.108}
\end{equation}
因此, 把 \eqref{eq:7.106} 与 $|x_3|\leq h$ 结合, 椭圆理论 (例如
\cite{GiTr} 推论 6.3) 给出
\begin{equation}
|x|\bigl|\nabla_{\Sigma}x_3\bigr|
+|x|^2\bigl|\nabla_{\Sigma}^2x_3\bigr|
\leq Ch+O\bigl(|x|^{-1}\bigr).
\label{eq:7.109}
\end{equation}
对梯度有界的图, $x_3$ 的 Hessian 与第二基本形式本质上是等价的 (二者
互相上下界), 故可断定
\begin{equation}
|x|^4|A(x)|^2\leq C'.
\label{eq:7.110}
\end{equation}

估计 \eqref{eq:7.105}--\eqref{eq:7.110} 使我们能抽取收敛子列与稳定的
极限 $\Sigma$, 它同样满足 \eqref{eq:7.105}、\eqref{eq:7.109}、
\eqref{eq:7.110} 以及环绕条件 $\Sigma\cap\sigma\neq\emptyset$. 由于每个
$\Sigma_s$ 都是嵌入的, 由强极值原理 $\Sigma$ 也是嵌入的. 把
\eqref{eq:7.105} 与 \eqref{eq:7.109} 结合, 得到
\begin{equation}
\int_{\Sigma}|A|^2<\infty.
\label{eq:7.111}
\end{equation}
再把它与 \eqref{eq:7.96} 及 Gauss 方程结合, 又有
\begin{equation}
\int_{\Sigma}|K_{\Sigma}|<\infty.
\label{eq:7.112}
\end{equation}
由于 $\Sigma$ 具有二次面积增长 (即 \eqref{eq:7.105}), 我们得到 (对数型)
截断函数 $0\leq\phi_j\leq1$ 满足
\begin{equation}
E(\phi_j)\leq1/j\to0\quad\text{且}\quad
\phi_j=1\quad\text{于}\ \{|x|\leq j\}\cap\Sigma.
\label{eq:7.113}
\end{equation}
用 $\phi_j$ 代入稳定性不等式, 给出
\begin{equation}
\int_{\Sigma}\bigl(|A|^2/2+\operatorname{Scal}_M-K_{\Sigma}\bigr)\phi_j^2
\leq E(\phi_j)\leq1/j.
\label{eq:7.114}
\end{equation}
利用 \eqref{eq:7.111}、\eqref{eq:7.112} 并令 $j\to\infty$, 由
\eqref{eq:7.114} 得到
\begin{equation}
0<\int_{\Sigma}\bigl(|A|^2/2+\operatorname{Scal}_M\bigr)
\leq\int_{\Sigma}K_{\Sigma}<\infty.
\label{eq:7.115}
\end{equation}

剩下要导出矛盾. 关键在于 $\mathcal{M}<0$ 蕴含 $\Sigma$ ``渐近地大于''
$\R^2$, 这与 \eqref{eq:7.115} 矛盾. 我们用 Gauss--Bonnet 定理把这一点
说清楚.

对 $|x|$ 大, $\Sigma$ 的每个连通分支都是 $\{x_3=0\}$ 上的图, 故 $\Sigma$
在大球内的面积由分支的个数决定. 由 \eqref{eq:7.105}, 只有一个分支; 即
对 $|x|$ 大, $\Sigma$ 是 $\{x_3=0\}$ 上某个函数 $u$ 的单值图. 由
\eqref{eq:7.109},
\begin{equation}
|x|^2\bigl|\nabla^2u\bigr|+|x||\nabla u|+u=O(1).
\label{eq:7.116}
\end{equation}
置 $\sigma_s=\{x_1^2+x_2^2=s^2\}\cap\Sigma
=\{(s\cos\theta,s\sin\theta,u(s\cos\theta,s\sin\theta)\}$. 利用
\eqref{eq:7.96} 与 \eqref{eq:7.116}, $\sigma_s$ 的测地曲率 (在 $M$ 中)
为 $s^{-1}+O(s^{-2})$. 再结合 \eqref{eq:7.110}, $\sigma_s$ 的测地曲率
$k_g$ (在 $\Sigma$ 中) 也是 $s^{-1}+O(s^{-2})$, 从而
\begin{equation}
\int_{\sigma_s}k_g
=\bigl(2\pi s+O(1)\bigr)\bigl(s^{-1}+O(s^{-2})\bigr)
=2\pi+O\bigl(s^{-1}\bigr).
\label{eq:7.117}
\end{equation}
另一方面, 对 $s$ 大, Gauss--Bonnet 定理给出
\begin{equation}
\int_{\sigma_s}k_g=2\pi-\int_{\{|x|<s\}\cap\Sigma}K.
\label{eq:7.118}
\end{equation}
当 $s$ 充分大时, \eqref{eq:7.115} 与 \eqref{eq:7.118} 同 \eqref{eq:7.117}
矛盾, 证毕.
\end{proof}

% ================================================================
%  7.7  Ricci 流的熄灭
% ================================================================
\section{Ricci 流的熄灭}
\label{sec:7.7}

本节证明: 在任意同伦三维球面上, Ricci 流在有限时间内熄灭. 这一结果属于
Colding--Minicozzi \cite{CM19} 与 Perelman \cite{Pe1}. 这里的证明遵循
\cite{CM19}. 证明的想法是估计宽度 (见 5.1 小节中的定义) 的变化率, 并证明
宽度在有限时间内变为零, 从而给出熄灭. 定理~\ref{thm:5.26} 证明了好的扫掠的存在性,
它在估计宽度的变化率时将起关键作用.

设 $M^3$ 是光滑闭可定向三流形, $g(t)$ 是 $M$ 上按 Hamilton Ricci 流
\cite{Ha1} 演化的一族度量, 即
\begin{equation}
\partial_tg=-2\operatorname{Ric}_{M_t}.
\label{eq:7.119}
\end{equation}
称三流形 $M$ 为\emph{素的}, 若它不能以非平凡的方式写成一个连通和 (与
$\bS^3$ 的连通和是平凡的). 若 $M$ 的所有高阶同伦群都消失, 则称 $M$ 为
\newterm{aspherical}{无球的}{aspherical}; 否则称 $M$ 为非无球的. 最后, 若
$M$ 与 $\bS^3$ 有相同的同伦群, 则称 $M$ 为一个同伦三维球面.

当 $M$ 素且非无球时, 由标准拓扑可知 $\pi_3(M)$ 非平凡 (例如可参看
\cite{CM19}). 对这种 $M$, 固定一个非平凡的同伦类 $\beta\in\Omega$. 于是
对每个度量 $g(t)$, 宽度
\[
W(g(t))=W(\beta,g(t))
\]
都是正的. 这一正性是下面关于 Ricci 流下宽度导数上界的定理中, 唯一用到
$M$ 的拓扑假设的地方. 作为推论, 我们得到该流在有限时间内熄灭 (手术发生
时熄灭时间的精确定义可参看 \cite{Pe1} 第 4.4 段).

\begin{theorem}[Colding--Minicozzi, \cite{CM19}]
\label{thm:7.19}
设 $M^3$ 是闭可定向、素且非无球的三流形, 并带有度量 $g=g(0)$. 在 Ricci
流之下, 宽度 $W(g(t))$ 满足
\begin{equation}
\frac{d}{dt}W(g(t))\leq-4\pi+\frac3{4(t+C)}W(g(t)),
\label{eq:7.120}
\end{equation}
其含义是前向差商的上极限. 于是 $g(t)$ 在有限时间内熄灭.
\end{theorem}

\eqref{eq:7.120} 中的 $4\pi$ 来自 Gauss--Bonnet 定理, 而 $3/4$ 来自演化
方程所蕴含的标量曲率最小值的界. 这两个常数都是要紧的; 而常数 $C>0$
依赖于初始度量, 其具体取值并不重要.

为看出 \eqref{eq:7.120} 蕴含有限的熄灭时间, 把 \eqref{eq:7.120} 改写为
\begin{equation}
\frac{d}{dt}\Bigl(W(g(t))(t+C)^{-3/4}\Bigr)\leq-4\pi(t+C)^{-3/4}
\label{eq:7.121}
\end{equation}
并积分, 得到
\begin{equation}
(T+C)^{-3/4}W(g(T))\leq C^{-3/4}W(g(0))
-16\pi\bigl[(T+C)^{1/4}-C^{1/4}\bigr].
\label{eq:7.122}
\end{equation}
由于按定义 $W\geq0$, 而 \eqref{eq:7.122} 的右端当 $T$ 充分大时变为负值,
即得所要结论.

特别地, 定理~\ref{thm:7.19} 表明: 对任意同伦三维球面, Ricci 流都会熄灭.
事实上, 作为推论我们还得到: 在不含无球和项的一切三流形上, Ricci 流都在
有限时间内熄灭 (这为何容易推出, 可参看 \cite{Pe1} 的 1.5 或 \cite{CM19}
第 4 节):

\begin{corollary}[Colding--Minicozzi, \cite{CM19}; Perelman, \cite{Pe1}]
\label{cor:7.20}
设 $M^3$ 是闭可定向三流形, 其素分解只含非无球的因子, 并带有度量
$g=g(0)$. 在带手术的 Ricci 流之下, $g(t)$ 在有限时间内熄灭.
\end{corollary}

Perelman 之所以关心有限时间熄灭的问题, 部分原因如下: 若有人对同伦三维
球面 (或更一般地, 不含无球和项的三流形) 的几何化感兴趣, 并且已经知道
Ricci 流在有限时间内熄灭, 那他就不必分析时间趋于无穷时流的行为. 于是
特别地, 也就不必使用塌缩论证.

证明定理~\ref{thm:7.19} 的关键要素之一, 是 $M$ 的一列好的扫掠的存在性:
该扫掠中面积接近宽度 (即扫掠中任一映射的最大能量) 的每个映射, 自身必须
接近一族调和映射. 这正是定理~\ref{thm:5.26} 所给出的.

我们将使用两个关于 $\bS^2$ 到流形的映射的接近与收敛的概念. 第一个 (也
较弱) 是第 3 章 \eqref{eq:3.36} 中定义的变分流形距离 $d_V$. 第二个是
第 3 章末节中定义的气泡收敛; 命题~\ref{prop:3.27} 表明气泡收敛蕴含变分流形收敛.

\mtpara{宽度变化率的上界}
本小节自始至终设 $M^3$ 是光滑闭、素且非无球的可定向三流形, 而 $g(t)$
是 $M$ 上按 Ricci 流演化的一族度量. 我们将证明定理~\ref{thm:7.19}, 给出
Ricci 流下宽度 $W(g(t))$ 之导数的上界. 为此需要三件事.

其一是标量曲率 $R=R(t)$ 的演化方程 (参看 \cite{Ha2} 第 16 页)
\begin{equation}
\partial_tR=\Delta R+2|\operatorname{Ric}|^2
\geq\Delta R+\frac23R^2,
\label{eq:7.123}
\end{equation}
它经一个直接的最大值原理论证蕴含: 在时刻 $t>0$,
\begin{equation}
R(t)\geq\frac1{1/[\min R(0)]-2t/3}=-\frac3{2(t+C)}.
\label{eq:7.124}
\end{equation}
曲率已作归一化, 使得在单位 $\bS^3$ 上 Ricci 曲率为 $2$, 标量曲率为 $6$.
在 \eqref{eq:7.124} 的推导中我们隐含假定 $\min R(0)<0$. 若并非如此, 则
\eqref{eq:7.124} 对任意 $C>0$ 都平凡地成立, 因为由 \eqref{eq:7.123},
$\min R(t)$ 总是单调不减的. 这最后一点在手术发生时也会用到, 这是因为
按构造, 任何手术区域的标量曲率都很大 (且为正).

证明中需要的第二件事是一个观察: 若 $\{\Sigma_i\}$ 是一族分支极小
$2$-球面, $f\in W^{1,2}(\bS^2,M)$ 且
$d_V\bigl(f,\bigcup_i\Sigma_i\bigr)<\epsilon$, 则对 $M$ 上任意光滑二次型
$Q$ 有 (单位法向量 $N_f$ 在 $J_f\neq0$ 处有定义)
\begin{equation}
\begin{aligned}
&\Bigl|\int_f\bigl[\operatorname{Tr}(Q)-Q(N_f,N_f)\bigr]
-\sum_i\int_{\Sigma_i}\bigl[\operatorname{Tr}(Q)
-Q(N_{\Sigma_i},N_{\Sigma_i})\bigr]\Bigr|\\
&\qquad<C\epsilon\|Q\|_{C^1}\Area(f).
\end{aligned}
\label{eq:7.125}
\end{equation}

最后一件事是极小 $2$-球面面积变化率的一个上界. 设 $X$ 是闭曲面,
$f:X\to M$ 是 $W^{1,2}$ 映射, 则利用 \eqref{eq:7.119} 容易算得 (参看
\cite{Ha2} 第 38--41 页)
\begin{equation}
\frac{d}{dt}\Big|_{t=0}\Area_{g(t)}(f)
=-\int_f\bigl[R-\operatorname{Ric}_M(N_f,N_f)\bigr].
\label{eq:7.126}
\end{equation}
若 $\Sigma\subset M$ 是闭浸入极小曲面, 则
\begin{equation}
\frac{d}{dt}\Big|_{t=0}\Area_{g(t)}(\Sigma)
=-\int_{\Sigma}K_{\Sigma}
-\frac12\int_{\Sigma}\bigl[|A|^2+R\bigr].
\label{eq:7.127}
\end{equation}
这里 $K_{\Sigma}$ 是 $\Sigma$ 的 (内蕴) 曲率. 为从 \eqref{eq:7.126} 得到
\eqref{eq:7.127}, 我们用到了: 若 $K_M$ 是 $M$ 在与 $\Sigma$ 相切的二平面
上的截面曲率, 则由 Gauss 方程与 $\Sigma$ 的极小性,
$K_{\Sigma}=K_M-\frac12|A|^2$. 下一个引理给出上述上界.

\begin{lemma}
\label{lem:7.21}
若 $\Sigma\subset M^3$ 是 $2$-球面的分支极小浸入, 则
\begin{equation}
\frac{d}{dt}\Big|_{t=0}\Area_{g(t)}(\Sigma)
\leq-4\pi-\frac{\Area_{g(0)}(\Sigma)}2\min_MR(0).
\label{eq:7.128}
\end{equation}
\end{lemma}

\begin{proof}
设 $\{p_i\}$ 是 $\Sigma$ 的分支点集, $b_i>0$ 为分支的阶. 由
\eqref{eq:7.127},
\begin{equation}
\frac{d}{dt}\Big|_{t=0}\Area_{g(t)}(\Sigma)
\leq-\int_{\Sigma}K_{\Sigma}-\frac12\int_{\Sigma}R
=-4\pi-2\pi\sum b_i-\frac12\int_{\Sigma}R,
\label{eq:7.129}
\end{equation}
其中等号用了带分支点的 Gauss--Bonnet 定理 (这一等式也可由 \cite{ScYa2}
第 10 页关于曲面间调和映射的 Bochner 型公式, 以及 \cite{ScYa2} 第 12 页
计入分支点的第二个显示式得到). 注意分支点只会对不等式 \eqref{eq:7.128}
有利.
\end{proof}

利用这三件事, 就可以证明宽度的变化率的上界.

\begin{proof}[定理~\ref{thm:7.19} 的证明]
固定一个时刻 $\tau$. 以下 $\tilde{C}$ 表示只依赖于 $\tau$ 的常数, 但允许
逐行改变. 设 $\gamma^j(\tau)$ 是定理~\ref{thm:5.26} 给出的、相应于度量 $g(\tau)$
的一列扫掠. 我们将用时刻 $\tau$ 的扫掠作为比较, 得到 $t>\tau$ 时宽度的
上界. 关键在下述断言: 给定 $\epsilon>0$, 存在 $\bar{j}$ 与 $\bar{h}>0$,
使得当 $j>\bar{j}$ 且 $0<h<\bar{h}$ 时有
\begin{equation}
\begin{aligned}
&\Area_{g(\tau+h)}\bigl(\gamma_s^j(\tau)\bigr)
-\max_{s_0}\Area_{g(\tau)}\bigl(\gamma_{s_0}^j(\tau)\bigr)\\
&\qquad\leq\Bigl[-4\pi+\tilde{C}\epsilon
+\frac3{4(\tau+C)}\max_{s_0}\Area_{g(\tau)}
\bigl(\gamma_{s_0}^j(\tau)\bigr)\Bigr]h+\tilde{C}h^2.
\end{aligned}
\label{eq:7.130}
\end{equation}
为看出 \eqref{eq:7.130} 为何蕴含 \eqref{eq:7.120}, 利用宽度两种定义
的等价性得到
\begin{equation}
W(g(\tau+h))
\leq\max_{s\in[0,1]}\Area_{g(\tau+h)}\bigl(\gamma_s^j(\tau)\bigr),
\label{eq:7.131}
\end{equation}
再在 \eqref{eq:7.130} 中令 $j\to\infty$ (于是
$\max_{s_0}\Area_{g(\tau)}\bigl(\gamma_{s_0}^j(\tau)\bigr)\to W(g(\tau))$),
得到
\begin{equation}
\frac{W(g(\tau+h))-W(g(\tau))}h
\leq-4\pi+\tilde{C}\epsilon
+\frac3{4(\tau+C)}W(g(\tau))+\tilde{C}h.
\label{eq:7.132}
\end{equation}
在 \eqref{eq:7.132} 中令 $\epsilon\to0$, 即得 \eqref{eq:7.120}.

剩下要证 \eqref{eq:7.130}. 首先, 设 $\delta>0$, 而 $\bar{j}$ 由定理~\ref{thm:5.26}
给出, 它依赖于 $\epsilon$ (以及 $\tau$). 若 $j>\bar{j}$ 且
\[
\Area_{g(\tau)}\bigl(\gamma_s^j(\tau)\bigr)>W(g)-\delta,
\]
则令 $\bigcup_i\Sigma_{s,i}^j(\tau)$ 为定理~\ref{thm:5.26} 给出的极小球面族. 把
\eqref{eq:7.126}、取 $Q=\operatorname{Ric}_M$ 的 \eqref{eq:7.125} 以及
引理~\ref{lem:7.21} 结合, 得到
\begin{equation}
\begin{aligned}
\frac{d}{dt}\Big|_{t=\tau}\Area_{g(t)}\bigl(\gamma_s^j(\tau)\bigr)
&\leq\frac{d}{dt}\Big|_{t=\tau}\Area_{g(t)}
\Bigl(\bigcup_i\Sigma_{s,i}^j(\tau)\Bigr)\\
&\quad+\tilde{C}\epsilon\|\operatorname{Ric}_M\|_{C^1}
\Area_{g(\tau)}\bigl(\gamma_s^j(\tau)\bigr)\\
&\leq-4\pi
-\frac{\Area_{g(\tau)}\bigl(\gamma_s^j(\tau)\bigr)}2\min_MR(\tau)
+\tilde{C}\epsilon.
\end{aligned}
\label{eq:7.133}
\end{equation}
代入 $R(\tau)$ 的下界 \eqref{eq:7.124}, 给出
\begin{equation}
\frac{d}{dt}\Big|_{t=\tau}\Area_{g(t)}\bigl(\gamma_s^j(\tau)\bigr)
\leq-4\pi
+\frac{3\max_{s_0}\Area_{g(\tau)}\bigl(\gamma_{s_0}^j(\tau)\bigr)}
{4(\tau+C)}+\tilde{C}\epsilon.
\label{eq:7.134}
\end{equation}
由于度量 $g(t)$ 光滑变化, 且每个扫掠 $\gamma^j$ 的能量一致有界, 容易
看出 $\Area_{g(\tau+h)}\bigl(\gamma_s^j(\tau)\bigr)$ 是 $h$ 的光滑函数,
且在 $h=0$ 附近有一致于 $j$ 与 $s$ 的 $C^2$ 界 (参看 \eqref{eq:7.126}).
特别地, \eqref{eq:7.134} 与 Taylor 展开给出 $\bar{h}>0$ (不依赖于 $j$),
使得对满足
\[
\Area_{g(\tau)}\bigl(\gamma_s^j(\tau)\bigr)>W(g)-\delta
\]
的 $s$, \eqref{eq:7.130} 成立.

在剩下的情形有
\[
\Area\bigl(\gamma_s^j(\tau)\bigr)\leq W(g)-\delta,
\]
于是 $g(t)$ 的连续性蕴含 (必要时再缩小 $\bar{h}>0$) \eqref{eq:7.130}
自动成立.
\end{proof}

\mtpara{手术发生时}
发生手术时有两个要点. 第一, 标量曲率在手术区域为正, 因此标量曲率的下界
$R\geq-\frac3{2(t+C)}$ 在出现手术时仍然有效. 第二, 若存在手术, 不等式
\eqref{eq:7.122} 仍然有效. 具体地说, 设 $M$ 如推论~\ref{cor:7.20} 所述,
并且是不可约的. 由一次手术可以构造一个从手术前流形到手术后流形某个分支
的度为 $1$ 的映射 (把无关的``几何''分支丢掉). 用 Perelman 的手术方法
(参看 \cite{Pe2}), 该映射是 $(1+\xi)$-Lipschitz 的, 其中函数 $\xi$ 在
\cite{KL} 的注 73.6 中有描述. 函数 $\xi(t)$ 依赖于手术前流形所处的时刻
$t$, 并在 $t$ 趋于手术时刻时趋于零.
